% !PRINTED 206
\weberpage{206}
\begin{center}
\S.\,67.

\gesperrt{Das volle Invariantensystem der binären\\
biquadratischen Form.}
\end{center}

Wir wollen noch beweisen, dass mit den Bildungen $A$, $B$ das System der
unabhängigen Invarianten der binären, biquadratischen Form erschöpft ist.
Während man sich aber gewöhnlich mit dem Nachweis begnügt, dass jede
Invariante eine ganze rationale Function von $A$, $B$ ist, wollen wir, wie wir
es schon in \S.\,63 für die Covarianten der cubischen Form gethan haben, auch
die Frage nach den numerischen Co\"efficienten berühren, und hier zeigt sich
eine neue Erscheinung.

Die Invariante $D$ z.\,B. ist zwar nach \S.\,64, (11) rational durch $A$, $B$
ausgedrückt. Aber die Zahlenco\"efficienten sind nicht ganze Zahlen, sondern
haben den Nenner 27, obwohl die Co\"efficienten in der entwickelten Function
$D$ alle ganze Zahlen sind.

Wir betrachten also jetzt als ganze Invarianten ganze rationale homogene
Functionen $\mu^{\mathrm{ten}}$ Grades der fünf Veränderlichen $a_0, a_1, a_2, a_3, a_4$
\weq{1}{
I(a_0, a_1, a_2, a_3, a_4) = I(a),
}
deren Co\"efficienten ganze Zahlen sind, und denen die Invarianteneigenschaft
\weq{2}{
I' = r^{2\mu} I
}
zukommt, wenn $I'$ oder $I(a')$ dieselbe Function der Co\"efficienten einer
transformirten Function ist.

Solche Invarianten sind $A$, $B$, $D$, zwischen denen die Relation besteht
\weq{3}{
27 D = 4A^3 - B^2,
}
und unser Ziel ist, zu beweisen, dass alle ganzen Invarianten ganze und
ganzzahlige rationale Functionen von diesen dreien sind.

Wenn wir die Function $I'$ in der Normalform (4) des \S.\,65 bilden, so
erhalten wir zunächst aus (2), da $r = 1$ ist,
\[
I(a) = I'(0, a_1', a_2', a_3', 0),
\]
und diese Function $I$ kann nur von $a_2'$ und dem Product $a_1'a_3'$
abhängen; denn die Substitution
%CONT>>
